\section*{Lab 4 (DBSCAN) Implementation Tasks}
\textbf{Lab 4 focus:} density-based clustering, outlier detection, scaling sensitivity, and DBSCAN behavior on toy and Iris datasets.

\medskip
\textbf{Implementation policy (strict):}
\begin{itemize}
  \item Implement the DBSCAN core algorithm with NumPy only.
  \item Do not use a library DBSCAN implementation for the main task.
  \item Plotting libraries are allowed for visual inspection.
  \item Work through the notebook in order.
  \item Use \texttt{min\_pts} in code, which corresponds to \text{MinPts} in the lecture.
\end{itemize}

\section*{Core Equations And Rules}
\begin{itemize}
  \item Dataset:
  \[
    X = \{x_1,\dots,x_n\}
  \]
  \item Epsilon-neighborhood:
  \[
    N_\varepsilon(x_i) = \{x_j : \|x_j - x_i\|_2 \le \varepsilon\}
  \]
  \item Core, border, and noise points:
  \[
    \text{core if } |N_\varepsilon(x_i)| \ge \text{MinPts}
  \]
  \[
    \text{border if it lies in a core neighborhood but is not core}
  \]
  \[
    \text{noise if it is neither core nor border}
  \]
\end{itemize}

\section*{Data Files And Preparation}
\begin{itemize}
  \item \texttt{dbscan.ipynb}: guided notebook with the assignment steps.
  \item \texttt{advanced\_ai/clustering/dbscan.py}: starter implementation.
  \item \texttt{prepare\_datasets.py}: generates the prepared data bundles.
  \item \texttt{scripts/prepare\_all\_data.sh}: shell entrypoint for dataset preparation.
  \item \texttt{scripts/check\_prepare\_datasets.sh}: verification script for the prepared bundles.
  \item \texttt{iris\_dbscan.py}, \texttt{toy\_dbscan.py}, \texttt{scaling\_and\_outliers.py}: task runners.
\end{itemize}

\medskip
Expected prepared files:
\begin{itemize}
  \item \texttt{datasets/toy\_dbscan.npz}
  \item \texttt{datasets/iris\_dbscan.npz}
  \item \texttt{datasets/scaling\_study.npz}
  \item \texttt{datasets/outlier\_study.npz}
\end{itemize}

\medskip
The scripts use the conda environment \texttt{ai\_courses}.

\section*{Roadmap}
\begin{itemize}
  \item Start with Step 1.
  \item After that, continue to the next section in order.
  \item Each step introduces only the task you need next.
\end{itemize}

\section*{Step 1: Implement the DBSCAN Core}
\begin{itemize}
  \item Open \texttt{advanced\_ai/clustering/dbscan.py} and fill in the missing methods.
  \item Start with the neighborhood query and cluster expansion, then complete \texttt{fit} and \texttt{fit\_predict}.
\end{itemize}

\section*{Step 2: Run the Self-Checks}
\begin{itemize}
  \item After implementing the core class, run the self-check suite and make sure every check passes before continuing.
\end{itemize}

\section*{Step 3: Warm-up by Hand}
\begin{itemize}
  \item Take the 2D toy dataset from the notebook and classify the marked points as core, border, or noise using the given $\varepsilon$ and \text{MinPts}.
  \item Write a short justification for each answer.
\end{itemize}

\section*{Step 4: Toy Dataset}
\begin{itemize}
  \item Run \texttt{toy\_dbscan.py} on the prepared 2D dataset.
  \item Report the number of clusters and noise points.
  \item Save the cluster plot.
  \item Compare the result with the hand-worked example.
\end{itemize}

\section*{Step 5: Iris Data}
\begin{itemize}
  \item Run \texttt{iris\_dbscan.py} on the prepared Iris bundle.
  \item Compare raw and standardized features.
  \item Inspect how $\varepsilon$ changes the result.
  \item Record cluster sizes and the noise fraction.
\end{itemize}

\section*{Step 6: Scaling and Outliers}
\begin{itemize}
  \item Run \texttt{scaling\_and\_outliers.py}.
  \item Show why feature scaling matters for DBSCAN.
  \item Identify which points are noise in the prepared outlier dataset.
  \item Explain why these points are not border points.
\end{itemize}

\section*{Step 7: Optional Extension}
\begin{itemize}
  \item A small sweep over $\varepsilon$ and \text{MinPts}.
  \item A simple plot of the number of clusters versus $\varepsilon$.
  \item A short comparison of DBSCAN and $k$-means on the same dataset.
\end{itemize}
\section*{Submission}
\begin{itemize}
  \item Submit the notebook, the completed starter files, and the generated output files from the required runners.
\end{itemize}
